\chapter{Developer documentation}
\label{ch:impl}

This chapter provides comprehensive developer documentation for the code coverage visualization feature, including the solution plan, implementation details, and testing documentation.

\section{Solution Plan}

\subsection{System Architecture}

The code coverage feature extends CodeChecker's three-tier architecture:

\begin{itemize}
    \item \textbf{Data Layer}: Database storage for coverage information using SQLAlchemy ORM
    \item \textbf{Service Layer}: Thrift-based API endpoints for coverage data retrieval
    \item \textbf{Presentation Layer}: Vue.js web interface components for visualization
\end{itemize}

In addition, a new \texttt{report-converter} analyzer module handles the conversion of LCOV data into CodeChecker's internal format. This follows the same pattern used by all other external analyzer integrations in CodeChecker (e.g., Pylint, ESLint, Cppcheck).

\begin{figure}[H]
\centering
\begin{tabular}{|c|}
\hline
\textbf{User's Build System} \\
\texttt{gcc --coverage} $\rightarrow$ \texttt{./run\_tests} $\rightarrow$ \texttt{lcov --capture} \\
$\downarrow$ \texttt{coverage.info} \\
\hline
\hline
\textbf{Report Converter} (\texttt{report-converter -t lcov}) \\
$\downarrow$ \texttt{coverage/coverage.json} \\
\hline
\hline
\textbf{CodeChecker Store} (\texttt{CodeChecker store}) \\
$\downarrow$ \\
\hline
\hline
\textbf{Database Layer} \\
\texttt{test\_coverage} + \texttt{test\_coverage\_summary} \\
$\downarrow$ \\
\hline
\hline
\textbf{Thrift API Layer} \\
\texttt{getTestCoverage()} + \texttt{getFilesWithCoverage()} \\
$\downarrow$ \\
\hline
\hline
\textbf{Vue.js Frontend} \\
\texttt{CoverageStatistics} + \texttt{CodeCoverageView} \\
\hline
\end{tabular}
\caption{Coverage feature data flow: from the user's build system through the converter, database, API, and frontend layers.}
\label{fig:architecture}
\end{figure}

\subsection{Technology Stack}

\begin{itemize}
    \item \textbf{Backend}: Python 3.x with SQLAlchemy~\cite{sqlalchemy} ORM for database operations
    \item \textbf{API}: Apache Thrift~\cite{thrift} for service definition and communication
    \item \textbf{Database}: PostgreSQL or SQLite (via CodeChecker's database abstraction)
    \item \textbf{Frontend}: Vue.js 2.x~\cite{vuejs} with Vuetify and CodeMirror~\cite{codemirror}
    \item \textbf{Coverage Format}: LCOV~\cite{lcov} \texttt{.info} files as input, JSON as intermediate format
\end{itemize}

\subsection{Database Design}

The coverage feature extends CodeChecker's database schema with two tables.

\begin{figure}[H]
\centering
\begin{tabular}{|l|l|l|}
\hline
\multicolumn{3}{|c|}{\textbf{files} (existing CodeChecker table)} \\
\hline
\texttt{id} & Integer & Primary key \\
\texttt{filepath} & String & Absolute file path \\
\texttt{content\_hash} & String & SHA hash of content \\
\hline
\end{tabular}

\vspace{0.5cm}

\begin{tabular}{|l|l|l|}
\hline
\multicolumn{3}{|c|}{\textbf{test\_coverage} (new)} \\
\hline
\texttt{id} & Integer & Primary key \\
\texttt{file\_id} & Integer & FK $\rightarrow$ \texttt{files.id}, indexed \\
\texttt{covered\_lines} & Integer & $+N$ = covered, $-N$ = uncovered \\
\hline
\end{tabular}

\vspace{0.5cm}

\begin{tabular}{|l|l|l|}
\hline
\multicolumn{3}{|c|}{\textbf{test\_coverage\_summary} (new)} \\
\hline
\texttt{id} & Integer & Primary key \\
\texttt{file\_id} & Integer & FK $\rightarrow$ \texttt{files.id}, unique \\
\texttt{functions\_found} & Integer & From LCOV \texttt{FNF:} \\
\texttt{functions\_hit} & Integer & From LCOV \texttt{FNH:} \\
\hline
\end{tabular}
\caption{Database schema for coverage data. The \texttt{test\_coverage} table uses sign encoding: positive values represent covered lines, negative values represent uncovered executable lines. The \texttt{test\_coverage\_summary} table stores per-file function coverage statistics.}
\label{fig:dbschema}
\end{figure}

\subsubsection{TestCoverage Table}

The \texttt{test\_coverage} table stores line-level coverage data. Each row represents a single instrumentable line for a specific file. Covered lines are stored as positive integers; uncovered (but executable) lines are stored as negative integers. This encoding allows a single table to represent both states while the API can distinguish them by sign.

\begin{lstlisting}[language=Python, caption={TestCoverage database model}]
class TestCoverage(Base):
    __tablename__ = "test_coverage"
    id = Column(Integer, autoincrement=True, primary_key=True)
    file_id = Column(Integer, ForeignKey('files.id',
                    ondelete='CASCADE',
                    deferrable=True,
                    initially="DEFERRED"),
                    index=True)
    covered_lines = Column(Integer)
\end{lstlisting}

The index on \texttt{file\_id} optimizes the primary access pattern: retrieving all coverage data for a specific file. The deferred foreign key constraint allows bulk insertion during the store operation.

\subsubsection{TestCoverageSummary Table}

The \texttt{test\_coverage\_summary} table stores per-file function coverage statistics, extracted from the LCOV \texttt{FNF} (functions found) and \texttt{FNH} (functions hit) records:

\begin{lstlisting}[language=Python, caption={TestCoverageSummary database model}]
class TestCoverageSummary(Base):
    __tablename__ = "test_coverage_summary"
    id = Column(Integer, autoincrement=True, primary_key=True)
    file_id = Column(Integer, ForeignKey('files.id',
                    ondelete='CASCADE',
                    deferrable=True,
                    initially="DEFERRED"),
                    index=True, unique=True)
    functions_found = Column(Integer, default=0)
    functions_hit = Column(Integer, default=0)
\end{lstlisting}

The \texttt{unique=True} constraint on \texttt{file\_id} ensures one summary row per file. Updates to existing data use upsert semantics.

\subsection{Module Structure}

\begin{description}
    \item[\texttt{analyzers/lcov/analyzer\_result.py}] LCOV report converter. Parses \texttt{.info} files and outputs \texttt{coverage.json}.

    \item[\texttt{run\_db\_model.py}] Database model definitions: \texttt{TestCoverage} and \texttt{TestCoverageSummary}.

    \item[\texttt{mass\_store\_run.py}] Storage integration: \texttt{\_\_add\_coverage} method processes \texttt{coverage.json} during \texttt{CodeChecker store}.

    \item[\texttt{report\_server.py}] API endpoints: \texttt{getTestCoverage} and \texttt{getFilesWithCoverage}.

    \item[\texttt{CoverageStatistics.vue}] Directory-level coverage statistics dashboard.

    \item[\texttt{CodeCoverageView.vue}] File-level coverage visualization with line highlighting.

    \item[\texttt{report\_server.thrift}] Thrift API contract definitions.
\end{description}

\subsection{User Interface Design}

The web interface provides two views:

\begin{itemize}
    \item \textbf{Statistics view}: LCOV-style directory drill-down showing line and function coverage per directory/file, with coverage bars, percentages, and a search filter
    \item \textbf{File view}: Source code display with line-by-line green/red/white coloring based on LCOV data, coverage summary header, and searchable file selector
\end{itemize}

\section{LCOV Report Converter}

The LCOV converter is implemented as a new analyzer module in CodeChecker's \texttt{report-converter} tool, following the same pattern as all other converters. It is auto-discovered by the CLI through the dynamic module loading mechanism in \texttt{cli.py}.

\subsection{Converter Design}

The converter extends \texttt{AnalyzerResultBase} and overrides the \texttt{transform()} method instead of \texttt{get\_reports()}, because coverage data is not a traditional ``report'' with bug paths. This is the same approach used by the existing \texttt{gcov} converter.

\begin{lstlisting}[language=Python, caption={LCOV converter class structure}]
class AnalyzerResult(AnalyzerResultBase):
    TOOL_NAME = 'lcov'
    NAME = 'LCOV coverage info to JSON'
    URL = 'https://github.com/linux-test-project/lcov'

    def transform(self, analyzer_result_file_paths,
                  output_dir_path, export_type, ...):
        # Parse all .info files, merge results
        # Output coverage/coverage.json

    def get_reports(self, file_path):
        return []  # No Report objects for coverage
\end{lstlisting}

\subsection{LCOV Parsing}

The parser processes LCOV \texttt{.info} files line by line, extracting four types of records:

\begin{itemize}
    \item \texttt{SF:<path>} --- Start of a new source file section
    \item \texttt{DA:<line>,<count>} --- Line data with execution count
    \item \texttt{FNF:<n>} --- Number of functions found in the file
    \item \texttt{FNH:<n>} --- Number of functions hit (executed at least once)
    \item \texttt{end\_of\_record} --- End of the current file section
\end{itemize}

Unlike the original approach that only stored lines with \texttt{count > 0}, the converter now stores \emph{all} \texttt{DA:} records. Lines with \texttt{count > 0} are ``covered''; lines with \texttt{count = 0} are ``uncovered but executable.'' Lines without any \texttt{DA:} record are non-executable (comments, blank lines, declarations). This distinction is critical for accurate visualization---the compiler's instrumentation is the authoritative source for which lines are executable, not heuristic pattern matching on source code.

\begin{lstlisting}[language=Python, caption={LCOV parsing---extracting all DA records}]
elif line.startswith('DA:') and current_file:
    match = re.match(r'DA:(\d+),(\d+)', line)
    if match:
        line_num = int(match.group(1))
        exec_count = int(match.group(2))
        line_data[line_num] = exec_count  # Store ALL DA lines
\end{lstlisting}

\subsection{Output Format}

The converter produces \texttt{coverage/coverage.json} inside the output directory, structured so that \texttt{CodeChecker store} automatically detects it:

\begin{lstlisting}[language=Python, caption={Building the output JSON}]
for src, line_counts in coverage_map.items():
    covered = sorted(ln for ln, cnt in line_counts.items()
                     if cnt > 0)
    uncovered = sorted(ln for ln, cnt in line_counts.items()
                       if cnt == 0)
    entry = {"lines": covered, "uncovered": uncovered}
    if src in function_map:
        entry["functions_found"] = function_map[src]["found"]
        entry["functions_hit"] = function_map[src]["hit"]
    output[src] = entry
\end{lstlisting}

A key design choice in the parser is to process the \texttt{.info} file in a single pass, line by line, rather than using an XML or JSON parsing library. This is appropriate because the LCOV format is line-oriented and simple enough that a state machine (tracking the current source file) is sufficient. The alternative---loading the entire file into memory and using regular expressions to extract sections---would be less memory-efficient for large coverage files (the xerces-c \texttt{coverage.info} is 93,651 lines) and would not improve correctness.

The decision to store \emph{all} \texttt{DA:} records (both covered and uncovered) rather than only covered lines was made after discovering that the frontend's \texttt{isExecutableLine()} heuristic produced over 20 mismatches per file when compared against actual LCOV data. For example, the heuristic classified \texttt{return fUnicodeForm;} as non-executable (matching its ``simple return'' pattern), while the compiler instrumented it as an executable line. By storing the compiler's own determination of which lines are executable, the visualization becomes authoritative rather than approximate. The cost is roughly double the database rows (66,649 vs.\ 40,210 for covered-only), which is negligible for modern databases.

\section{Storage Integration}

Coverage data is automatically detected and stored during the \texttt{CodeChecker store} operation. The \texttt{MassStoreRun} class searches for \texttt{coverage/coverage.json} in the reports directory.

\subsection{File Path Matching}

A key challenge is matching file paths in \texttt{coverage.json} (absolute paths from the build machine) to file IDs in the database. The implementation uses six fallback strategies:

\begin{enumerate}
    \item Exact match against the current store operation's \texttt{file\_path\_to\_id} map
    \item Trimmed path match (removing configured path prefixes)
    \item Database query for exact path
    \item Database query for trimmed path
    \item Filename-only match against the current store
    \item Database query for filename-only match
\end{enumerate}

\subsection{Coverage Data Storage}

Covered lines are stored as positive integers and uncovered lines as negative integers in the \texttt{test\_coverage} table. Function coverage statistics are stored in \texttt{test\_coverage\_summary}:

\begin{lstlisting}[language=Python, caption={Storing coverage data with sign encoding}]
# Store covered lines (positive)
for line in new_lines:
    session.add(TestCoverage(file_id, line))

# Store uncovered lines (negative)
for uline in uncovered:
    if -uline not in existing_lines:
        session.add(TestCoverage(file_id, -uline))

# Store function summary
if fnf > 0:
    session.add(TestCoverageSummary(file_id, fnf, fnh))
\end{lstlisting}

The method also supports backward compatibility with the legacy format where \texttt{coverage.json} contains plain lists of covered line numbers.

The six-strategy path matching may appear over-engineered, but each strategy addresses a real-world scenario. Strategy 1 (exact match) handles the common case where the build machine and the CodeChecker server use the same paths. Strategy 2 (trimmed prefix) handles CI environments where build paths include job-specific prefixes (e.g., \texttt{/jenkins/workspace/job-123/src/file.c}). Strategies 3--4 (database queries) handle the case where coverage data is stored in a separate \texttt{CodeChecker store} invocation from the static analysis results, so the file is already in the database but not in the current store operation's path map. Strategies 5--6 (filename-only) are a last resort for cross-platform scenarios where the directory structure differs entirely. If all six strategies fail, the file is skipped with a warning---this is preferable to silently associating coverage data with the wrong file.

\section{API Endpoints}

\subsection{getTestCoverage}

Retrieves all coverage line numbers for a specific file. Returns both positive (covered) and negative (uncovered) values. The frontend uses the sign to determine coloring.

\begin{lstlisting}[language=Python, caption={getTestCoverage API implementation}]
def getTestCoverage(self, file_id):
    self.__require_view()
    with DBSession(self._Session) as session:
        q = session.query(TestCoverage.covered_lines) \
            .filter(TestCoverage.file_id == file_id)
        return [row[0] for row in q.all()]
\end{lstlisting}

\subsection{getFilesWithCoverage}

Returns all files with coverage data, including aggregated line counts and function statistics. This endpoint is optimized to return all data in a single query, avoiding the need for per-file API calls from the frontend:

\begin{lstlisting}[language=Python, caption={getFilesWithCoverage---single-query optimization}]
def getFilesWithCoverage(self):
    self.__require_view()
    with DBSession(self._Session) as session:
        # Covered line counts (positive values)
        covered_counts = dict(session.query(
            TestCoverage.file_id,
            func.count(TestCoverage.covered_lines)
        ).filter(TestCoverage.covered_lines > 0)
         .group_by(TestCoverage.file_id).all())

        # Total executable line counts (all DA lines)
        total_counts = dict(session.query(
            TestCoverage.file_id,
            func.count(TestCoverage.covered_lines)
        ).group_by(TestCoverage.file_id).all())

        # Function counts from summary table
        func_data = {row.file_id: (row.functions_found,
                     row.functions_hit)
                     for row in
                     session.query(TestCoverageSummary).all()}

        # Encode stats into SourceFileData fields
        for file_id, filepath in files:
            hit = covered_counts.get(file_id, 0)
            total = total_counts.get(file_id, 0)
            ff, fh = func_data.get(file_id, (0, 0))
            result.append(SourceFileData(
                fileId=file_id, filePath=filepath,
                fileContent=f"{hit}/{total}",
                remoteUrl=f"{ff}/{fh}" if ff > 0 else None))
\end{lstlisting}

The line count and function statistics are encoded into the \texttt{SourceFileData} response object's optional fields (\texttt{fileContent} for line counts as \texttt{"hit/total"}, \texttt{remoteUrl} for function stats as \texttt{"found/hit"}), avoiding the need to modify the Thrift API definition or regenerate client code.

Encoding data into existing fields is a pragmatic trade-off. The proper approach would be to define a new Thrift struct (e.g., \texttt{CoverageFileData} with explicit \texttt{lineHit}, \texttt{lineTotal}, \texttt{functionsFound}, \texttt{functionsHit} fields), regenerate the Python server stubs and JavaScript client code, and update all API consumers. This would require the Thrift compiler, which was not available in the development environment, and would affect the generated API packages distributed to all CodeChecker installations. The field-encoding approach achieves the same result with zero changes to the Thrift IDL, at the cost of a less self-documenting API. For a feature that is still being validated, this trade-off is appropriate---the encoding can be replaced with proper Thrift types once the feature is stable and the API contract is finalized.

\section{Frontend Implementation}

\subsection{CoverageStatistics Component}

The statistics page makes a single API call to \texttt{getFilesWithCoverage()} and builds the entire directory tree client-side from the file paths. Coverage numbers are parsed from the response's encoded fields:

\begin{lstlisting}[language=Java, caption={Parsing coverage data from API response}]
let lineHit = 0, lineTotal = 0;
if (file.fileContent) {
    const lp = file.fileContent.split("/");
    lineHit = parseInt(lp[0], 10) || 0;
    lineTotal = parseInt(lp[1], 10) || lineHit;
}
let fnFound = 0, fnHit = 0;
if (file.remoteUrl) {
    const parts = file.remoteUrl.split("/");
    fnFound = parseInt(parts[0], 10) || 0;
    fnHit = parseInt(parts[1], 10) || 0;
}
\end{lstlisting}

Directory-level aggregates are computed by summing the per-file values. The component includes a search filter (\texttt{v-text-field}) that filters the table by file or directory name.

\subsection{CodeCoverageView Component}

The file-level view fetches source content and coverage data for a single file. Line coloring uses the sign convention from the database:

\begin{lstlisting}[language=Java, caption={Sign-based line coloring}]
getLineCoverageStatus(lineNum) {
    const coveredSet = new Set(
        this.coveredLines.filter(l => l > 0));
    const uncoveredSet = new Set(
        this.coveredLines.filter(l => l < 0).map(l => -l));

    if (coveredSet.has(lineNum)) return "covered";   // green
    if (uncoveredSet.has(lineNum)) return "uncovered"; // red
    return "none";                                     // white
}
\end{lstlisting}

This approach eliminates the need for an \texttt{isExecutableLine()} heuristic that attempted to guess which lines are executable by pattern-matching source code. The compiler's instrumentation (via LCOV \texttt{DA:} records) is the authoritative source, and using it directly produces accurate results that match \texttt{gcov} and \texttt{lcov} HTML reports exactly.

The file selector uses Vuetify's \texttt{v-autocomplete} component with a custom filter function that searches both filename and full path.

The decision to build the directory tree client-side rather than server-side was driven by the single-query constraint. If the server returned a pre-built tree, it would need to know the user's current navigation state (which directory they are viewing), requiring a stateful API or additional query parameters. By returning a flat list of files with their coverage counts, the server remains stateless and the client can build any view (top-level directories, subdirectories, filtered results) without additional API calls. The scalability limit of this approach is the number of files: with 580 files, the JSON response is approximately 100~KB and the client-side tree construction takes under 50ms. For projects with tens of thousands of files, a server-side pagination or tree API would be necessary, but this is beyond the scope of the current implementation.

\section{Design Decisions}

This section describes the key design crossroads encountered during development and the reasoning behind each decision.

\subsection{Crossroad 1: Static Tool, Dynamic Data}

The most defining decision of the project concerned the scope of integration. The initial plan was ambitious: the user would provide only source code, and CodeChecker would handle the entire pipeline---compiling with coverage flags, executing the program, collecting \texttt{.gcda} files, running \texttt{lcov}, parsing the output, and storing the results. This would have made coverage a first-class feature, as simple as running \texttt{CodeChecker analyze}.

After extensive investigation, this approach proved fundamentally incompatible with CodeChecker's architecture. The core issue is that CodeChecker's workflow is built around the \emph{compilation database} (\texttt{compile\_commands.json}), which records how each source file is compiled but contains no information about how to \emph{execute} the resulting software. As my supervisor explained:

\begin{quote}
``The \texttt{CodeChecker analyze} command loops through the list of compilation actions and gives them to the static analyzers to read the source codes. However, a dynamic analyzer needs to \emph{run} the program. But in the compilation database there is no information about how to execute the software. Even worse, the compilation database doesn't even necessarily contain the command which produces an executable file. It is also possible that the project consists of multiple executables. How should an analyzer know which executables are to be run?''
\end{quote}

During the investigation, I discovered that some of the operations we wanted to automate (such as intercepting build commands to inject coverage flags) operate at the operating system level---the build logger works by injecting a shared library via \texttt{LD\_PRELOAD} that intercepts compiler invocations. Replicating this for test execution would require intercepting arbitrary program launches, which is far beyond the scope of a coverage tool.

The resolution was to follow the same pattern CodeChecker uses for all external analyzers: the user runs the tool themselves, and CodeChecker only handles the conversion and storage. This led to the \texttt{report-converter -t lcov} approach, where the user is responsible for the dynamic steps (compile, execute, capture) and CodeChecker takes over at the point where a \texttt{coverage.info} file already exists.

\subsection{Crossroad 2: Database Schema}

The second decision concerned how to store coverage data. Two approaches were considered:

\begin{description}
    \item[Option A: One row per file] Store the list of covered line numbers as a JSON column in a single database row per file. Simpler schema with fewer rows, and PostgreSQL's JSONB type supports efficient storage. However, this makes it impossible to query or aggregate individual lines at the SQL level, and extending with per-line metadata (such as execution counts) would require parsing the JSON in application code.

    \item[Option B: One row per line] Store each covered line as a separate row with a foreign key to the file. More rows, but enables SQL-level filtering, aggregation, and indexing. Follows CodeChecker's existing normalized database design.
\end{description}

Option B was chosen for consistency with CodeChecker's patterns and for future extensibility. The \texttt{test\_coverage} table stores one row per instrumentable line, with the sign encoding (positive = covered, negative = uncovered) added later to support accurate visualization without fetching source content.

The trade-off is storage volume: the xerces-c project produces 66,649 rows in \texttt{test\_coverage} plus 575 rows in \texttt{test\_coverage\_summary}. In practice, this is negligible for SQLite or PostgreSQL, and the indexed \texttt{file\_id} column keeps query times under 100ms even for the full dataset.

\subsection{Crossroad 3: User Experience Integration}

The third decision concerned where in CodeChecker's interface the coverage feature should live. Three options were considered:

\begin{description}
    \item[Option A: Extend \texttt{CodeChecker analyze}] Add coverage as another analyzer alongside Clang SA and Clang Tidy. This would have been the most seamless integration but was ruled out by the static-vs-dynamic incompatibility described above.

    \item[Option B: New \texttt{CodeChecker coverage} CLI command] A dedicated subcommand that wraps the entire coverage pipeline. This was implemented initially but removed because it only worked for trivial single-file projects and gave a false sense of completeness for real-world use.

    \item[Option C: Use \texttt{report-converter}] Treat LCOV as an external analyzer format, following the same pattern as Pylint, ESLint, and Cppcheck. The user runs their own coverage pipeline and converts the output.
\end{description}

Option C was chosen because it requires no changes to CodeChecker's core workflow, works with any project regardless of build system complexity, and follows the established pattern that CodeChecker developers and users already understand. The \texttt{CodeChecker coverage} command (Option B) was removed from the codebase to avoid maintaining a tool that would fail for most real projects.

\subsection{Why Sign Encoding for Uncovered Lines}

Storing uncovered lines as negative values in the same \texttt{test\_coverage} table avoids creating a second table and allows the existing \texttt{getTestCoverage} API to return both covered and uncovered lines in a single query. The frontend checks the sign to determine coloring: positive $\rightarrow$ green, negative $\rightarrow$ red, absent $\rightarrow$ white.

This encoding was introduced after discovering that the original frontend used an \texttt{isExecutableLine()} heuristic to guess which lines should be red. Comparing the heuristic's output against actual LCOV \texttt{DA:} records revealed over 20 mismatches per file---the compiler's instrumentation is the only reliable source for which lines are executable.

\subsection{Why Single-Query Statistics}

The statistics page initially made individual API calls for each file's coverage data and source content. With 580 files in the xerces-c test project, this produced 1,160 concurrent HTTP requests, causing \texttt{ERR\_INSUFFICIENT\_RESOURCES} in the browser. The solution was to return aggregated statistics from a single \texttt{getFilesWithCoverage()} call, encoding both hit count and total count into the \texttt{SourceFileData.fileContent} field as \texttt{"hit/total"}.

\section{Performance}

Performance was measured using the Apache Xerces-C++ project (580 source files, 40,210 covered lines, 26,439 uncovered lines, 66,649 total executable lines, 8,981 functions).

\begin{table}[H]
\centering
\begin{tabular}{|l|r|l|}
\hline
\textbf{Operation} & \textbf{Time} & \textbf{Notes} \\
\hline
LCOV conversion (\texttt{report-converter}) & $<$ 1 s & 93,651-line \texttt{.info} file $\rightarrow$ 414 KB JSON \\
\hline
\texttt{CodeChecker store} (with coverage) & $\sim$ 11 s & 580 files, 66,649 DB rows + 575 summaries \\
\hline
\texttt{getFilesWithCoverage} API call & $<$ 0.1 s & Single SQL query, returns all 580 files \\
\hline
Statistics page load (browser) & $<$ 1 s & Single API call, client-side aggregation \\
\hline
File-level view load & $<$ 0.5 s & Two API calls (source + coverage) \\
\hline
\end{tabular}
\caption{Performance measurements for the xerces-c project (580 files).}
\label{tab:performance}
\end{table}

The critical optimization was the single-query statistics page. The original design made two API calls per file (source content + coverage lines), producing 1,160 concurrent HTTP requests for 580 files. This caused \texttt{ERR\_INSUFFICIENT\_RESOURCES} in the browser. The fix---returning aggregated statistics from a single \texttt{getFilesWithCoverage()} call---reduced the page load from failure to under one second.

\section{Testing}

\subsection{Unit Tests}

The LCOV converter includes four unit tests:

\begin{description}
    \item[\texttt{test\_transform\_single\_info\_file}] Verifies correct parsing of a multi-file LCOV input, checking that covered line numbers match expected values.
    \item[\texttt{test\_transform\_empty\_file}] Verifies that an empty \texttt{.info} file produces no output.
    \item[\texttt{test\_get\_reports\_returns\_empty}] Confirms that coverage converters do not produce \texttt{Report} objects.
    \item[\texttt{test\_tool\_name}] Verifies the tool name used for CLI registration.
\end{description}

\subsection{End-to-End Verification}

The feature was tested end-to-end with the Apache Xerces-C++~\cite{xerces} project (a real-world XML parser with 80 test cases):

\begin{enumerate}
    \item Built xerces-c with \texttt{--coverage} flags (CMake)
    \item Ran all 80 tests via \texttt{ctest} (all passed)
    \item Captured coverage with \texttt{lcov} (335 \texttt{.gcda} files, 712 source files)
    \item Converted with \texttt{report-converter -t lcov} (580 files, 40,198 covered lines, 8,981 functions found, 5,952 functions hit)
    \item Stored in CodeChecker and verified via Thrift API and web UI
    \item Confirmed line coloring matches \texttt{lcov} HTML output
\end{enumerate}

\subsection{Lessons Learned}

\begin{description}
    \item[\textbf{isExecutableLine was wrong}] The original frontend used heuristic pattern matching to guess which lines are executable (e.g., skipping comments, closing braces, simple return statements). Comparison with actual LCOV \texttt{DA:} records showed 20+ mismatches per file. The fix was to use LCOV data directly.

    \item[\textbf{Browser resource limits}] Fetching source content for 580 files in parallel caused \texttt{ERR\_INSUFFICIENT\_RESOURCES}. The fix was server-side aggregation in a single SQL query.

    \item[\textbf{Lambda pickling in multiprocessing}] The server's \texttt{\_do\_db\_cleanups} function used a lambda with \texttt{ProcessPoolExecutor.map()}, which fails because lambdas cannot be pickled. The fix was to use \texttt{functools.partial} with \texttt{*zip(*products)}, matching the upstream CodeChecker code.

    \item[\textbf{CXX flags matter}] When building C++ projects with CMake, \texttt{CMAKE\_C\_FLAGS} only applies to C files. Omitting \texttt{CMAKE\_CXX\_FLAGS} results in only 2--3 \texttt{.gcda} files instead of 300+.
\end{description}
